% ============================================================================
% Two_Breaks — ApJ draft v2 (elevated, Burgess-informed).  aastex631.
% Single-pulse bright GRBs, empirical time-resolved spectroscopy:
%   (i) does Burgess+2014's thermal-non-thermal Ep-kT correlation appear in a
%       large sample?  (ii) what additional spectral properties/correlations emerge?
% %TODO (authors): finalize after the full-sample BB multi-start re-run + redshift
% cross-match.  %chk = double-check before submission.  Author/affil = placeholder.
% ============================================================================
\documentclass[twocolumn]{aastex631}

\newcommand{\nuc}{\ensuremath{\nu_c}}
\newcommand{\num}{\ensuremath{\nu_m}}
\newcommand{\Ep}{\ensuremath{E_{\rm p}}}
\newcommand{\kT}{\ensuremath{kT}}
\newcommand{\Liso}{\ensuremath{L_{\rm iso}}}
\newcommand{\daic}{\ensuremath{\Delta\mathrm{AIC}}}
\newcommand{\ssc}{\ensuremath{\sigma_{\rm sc}}}

\shorttitle{Empirical Spectral Properties of Single-Pulse GRBs}
\shortauthors{Khushboo et al.}

\begin{document}

\title{Time-Resolved Empirical Spectroscopy of a Large Sample of Single-Pulse
Gamma-Ray Bursts: a Large-Sample Test of the Thermal--Non-thermal Correlation
and a New Correlation Between the Two Spectral Breaks}

% %TODO author names, affiliations, ORCIDs are placeholders.
\author{Khushboo}
\affiliation{Affiliation TBD}
\author{Vikas Chand}
\affiliation{Affiliation TBD}
\email{vikas.chand.physics@gmail.com}
\author{Collaborators TBD}
\affiliation{Affiliation TBD}

\begin{abstract}
The curvature of gamma-ray burst (GRB) prompt spectra admits two physical
readings: a quasi-thermal photospheric (blackbody, BB) component on a
non-thermal continuum \citep{Burgess2014a}, or optically-thin synchrotron with
two break energies --- the cooling frequency \nuc\ and the injection frequency
$\num=\Ep$ \citep{Ravasio2019,Mei2024}. Using the shape-based selection of
\citet{Busby2024} we assemble a uniform sample of \textbf{106 single-pulse,
bright Fermi/GBM bursts} (18 with LAT/LLE) and fit \textbf{1063} Bayesian-block
time bins with a family of six \emph{empirical} models (Band, cutoff power law,
smoothly-broken power law, a double smoothly-broken power law [2SBPL], and
Band$+$BB / CPL$+$BB) by maximum likelihood. We deliberately use empirical
models --- model-agnostic, uniform across the sample, and not tuned to either
physical picture --- and ask whether the per-burst thermal--non-thermal
$\Ep$--$\kT$ correlation found by \citet{Burgess2014a} for six bursts appears at
scale, and what additional properties emerge. We find: (i) the low-energy photon
index peaks at $\alpha\simeq-0.75$, between the slow- ($-2/3$) and fast-cooled
($-3/2$) synchrotron limits and close to \citeauthor{Burgess2014a}'s $-0.81$;
(ii) of the bins requiring curvature beyond a single break, $78\%$ are equally or
better described by a thermal proxy (BB$+$SBPL) than by an explicit second break,
$94\%$ of all bins cannot decisively separate the six models, and the thermal
flux fraction is modest ($F_{\rm BB}/F_{\rm tot}\simeq0.1$); (iii) the
$\Ep$--$\kT$ correlation is recovered cleanly only in the highest signal-to-noise
bursts (e.g.\ GRB~130427A, $\rho=0.78$, $\Ep\propto\kT^{0.66}$), with empirical
$+$BB models detecting a blackbody in fewer bins than a physical synchrotron
model would; (iv) the two 2SBPL breaks are themselves correlated,
$\num\propto\nuc^{0.73\pm0.10}$ (Spearman $\rho=0.74$, intrinsic scatter
$\ssc=0.26$~dex, $N=40$) --- a relation not previously quantified
\citep{Ravasio2019,Mei2024}; and (v) in an observer-frame proxy for the
\citet{Mei2024} $\nuc$--\Liso\ relation, the apparent $\Ep$--brightness
correlation vanishes once an energy-weighting artifact is controlled (consistent
with their absence of an $\Ep$--\Liso\ relation), while a genuine
$\kT$--brightness correlation survives. We emphasize that we do \emph{not} fit a
physical synchrotron model; our results map the empirical observables and the
correlations among them, and we discuss their implications for jet composition.
\end{abstract}

\keywords{Gamma-ray bursts (629) --- Spectroscopy (1558) --- Non-thermal
radiation sources (1119) --- Astrostatistics (1882)}

% ============================================================================
\section{Introduction} \label{sec:intro}
% ============================================================================
GRB prompt-emission spectra are, to first order, universally described by the
empirical Band function \citep{Band1993} or a cutoff power law, characterized by
a low-energy photon index $\alpha$, a $\nu F_\nu$ peak energy \Ep, and (for Band)
a high-energy index $\beta$ \citep{Kaneko2006,Goldstein2012,Gruber2014,Yu2016}.
These functions fit the data but are physically agnostic: the same curvature can
arise from distinct mechanisms, and the inferred parameters do not by themselves
identify the radiative process \citep{KumarZhang2015}.

Two physically-motivated pictures dominate. In the first, part of the flux is
quasi-thermal emission from the jet photosphere, approximated by a blackbody (BB)
of temperature \kT, on a non-thermal continuum
\citep{MeszarosRees2000,Ryde2004,RydePeer2009,Peer2007,Guiriec2011}.
\citet{Burgess2014a}, fitting a \emph{physical} synchrotron$+$blackbody model to
six bright single-pulse bursts, reported a tight per-burst correlation
$\Ep\propto\kT^{\alpha}$, interpreting $\alpha\!\sim\!1$ as a baryon-dominated and
$\alpha\!\sim\!2$ as a magnetically-dominated jet (see also \citealt{Burgess2014b}).
In the second, the curvature is optically-thin synchrotron from a cooling
electron population, with \emph{two} break energies: the cooling frequency \nuc\
and the injection/peak frequency \num\ \citep{Sari1998}. A double smoothly-broken
power law (2SBPL) captures both breaks and has been used to argue for a
synchrotron origin \citep{Oganesyan2017,Oganesyan2018,Oganesyan2019,Ravasio2018,Ravasio2019}.
Most recently, \citet{Mei2024} showed that, under the synchrotron interpretation,
the fundamental spectral--energetics relation is $\nuc\propto\Liso^{0.53}$ rather
than the classical $\Ep$--\Liso\ \citep[Yonetoku;][]{Yonetoku2004} relation, which
is recovered only because the Band \Ep\ tracks \nuc.

Crucially, the two pictures produce similar curvature near the peak: a BB on a
single-break continuum (BB$+$SBPL) is an effective \emph{proxy} for a 2SBPL
(Figure~\ref{fig:schematic}). Distinguishing them, and characterizing the
resulting spectral properties at scale, is the goal of this work. Single-pulse
bursts are the cleanest laboratory: a single, non-overlapping pulse has the
simplest spectral evolution and avoids the blending of emission episodes
\citep{Burgess2014b,RydePeer2009,Busby2024}. \citet{Burgess2014a} were limited to
six bursts; whether their correlation is general, and what other properties a
large uniform sample reveals, has not been tested at scale.

We take a deliberately \emph{empirical} approach. Whereas \citet{Burgess2014a,Burgess2014b}
fit a physical synchrotron model, we fit a family of model-agnostic empirical
functions to a large, uniformly-selected sample (\S\ref{sec:sample}--\S\ref{sec:methods}),
and ask: (i) does the \citet{Burgess2014a} $\Ep$--$\kT$ correlation appear in a
large sample?, and (ii) what additional spectral correlations emerge --- in
particular between the two synchrotron breaks, and between break/temperature
energies and brightness? Results are in \S\ref{sec:results} and discussed in
\S\ref{sec:discussion}. Throughout we are explicit about the strengths (uniformity,
$N$, no model tuning) and the limits (an empirical continuum can absorb a real
component) of the empirical approach.

\begin{figure}[t]
\centering
\includegraphics[width=\columnwidth]{figures/fig_teach_schematic.png}
\caption{The two interpretations of the same spectral curvature. \emph{Left}: a
non-thermal continuum (Band) plus a thermal blackbody (\kT). \emph{Right}:
optically-thin synchrotron (2SBPL) with a cooling break \nuc\ and an
injection/peak break $\num=\Ep$. A BB$+$SBPL fit is a practical proxy for the 2SBPL.}
\label{fig:schematic}
\end{figure}

% ============================================================================
\section{Sample and Data} \label{sec:sample}
% ============================================================================
We select single-pulse bursts from the Fermi/GBM catalog using the shape-based
``horizontal-line'' algorithm of \citet{Busby2024}, which scores how well a
single FRED-like pulse describes the background-subtracted light curve. We retain
bursts above the \citeauthor{Busby2024} threshold (score $\gtrsim0.9983$) and
apply a fluence cut for adequate time-resolved signal-to-noise; no $T_{90}$ cut
is imposed. The selection is on pulse \emph{shape}, not brightness: the sample
fluence spans a factor of $\sim$250 (Figure~\ref{fig:sample}).

The final sample comprises \textbf{106 bursts} (Table~\ref{tab:sample}). For each
we use the Time-Tagged Event (TTE) data of all triggered NaI detectors within
$50^\circ$ of the source, the most-illuminated BGO, and, for the \textbf{18}
bursts with LAT Low-Energy (LLE) coverage, the LLE data above 30~MeV
\citep{Meegan2009,Atwood2009}. Joint fits use 8--900~keV for NaI (K-edge masked),
0.3--40~MeV for BGO, and 30--100~MeV for LLE.

\begin{figure}[t]
\centering
\includegraphics[width=\columnwidth]{figures/fig_teach_sample.png}
\caption{The single-pulse sample (\citealt{Busby2024} selection): fluence
distribution (left, $\sim$250$\times$ span) and $T_{90}$ versus fluence (right),
with the 18 LAT/LLE bursts highlighted.}
\label{fig:sample}
\end{figure}

% ============================================================================
\section{Methods} \label{sec:methods}
% ============================================================================
All spectral analysis uses the Multi-Mission Maximum Likelihood framework
\citep[3ML;][]{Vianello2015} with the \texttt{pgstat} likelihood (Poisson source,
Gaussian background).

\subsection{Background modeling} \label{sec:bkg}
GBM has no off-source pointing, so the background is large and time-variable. For
each detector we select pre- and post-burst intervals ($\sim$50--150~s wide each)
and fit a time polynomial (order chosen by AIC), interpolated beneath the burst.
Each detector is fit independently.

\subsection{Time binning: Bayesian blocks} \label{sec:bb}
Time bins are set by the Bayesian-blocks algorithm \citep{Scargle2013} on the
background-subtracted, 8--900~keV events of the brightest NaI detector, restricted
to the burst-emission interval (Figure~\ref{fig:bb}). We verified that our block
\emph{granularity} matches that of \citet{Burgess2014a} for the bursts in common
(typically within $\pm$2--4 blocks per burst); the residual edge differences arise
from the change-point prior and software (\citeauthor{Burgess2014b} used
\texttt{RMFIT} with prior $=8$ on the combined NaI$+$BGO events) and we show our
conclusions are robust to the binning (\S\ref{sec:nucnum}). The procedure yields
\textbf{1063} time bins.

\begin{figure}[t]
\centering
\includegraphics[width=\columnwidth]{figures/fig_teach_bayesian_blocks.png}
\caption{Example Bayesian-block binning (GRB~150902733, NaI, 8--900~keV). Red
segments are the blocks; dotted lines are the change-points defining the spectral
time bins.}
\label{fig:bb}
\end{figure}

\subsection{Spectral models} \label{sec:models}
In every bin we fit six empirical models. The Band function \citep{Band1993} is a
smoothly-joined broken power law: below the break
$E_b=(\alpha-\beta)\Ep/(2+\alpha)$,
\begin{equation}
N(E)=K\,(E/E_0)^{\alpha}\,e^{-(2+\alpha)E/\Ep},
\end{equation}
and above it $N(E)=K\,(E/E_0)^{\beta}\,[(\alpha-\beta)\Ep/(2+\alpha)E_0]^{\alpha-\beta}\,e^{\beta-\alpha}$,
with $E_0=100$~keV. The cutoff power law
(CPL) is $N(E)=K(E/E_0)^{\Gamma_{\rm ph}}\exp(-E/E_c)$, peaking at
$\Ep=(2+\Gamma_{\rm ph})E_c$. The smoothly-broken power law (SBPL) has indices
$\alpha,\beta$ and one break $E_b$; the double SBPL (2SBPL) has indices
$\alpha_1,\alpha_2,\beta$ and two breaks --- a lower break $x_b\equiv\nuc$ and an
upper break/peak $x_p\equiv\num=\Ep$. For both SBPL and 2SBPL the break-smoothness
(and the 2SBPL pivot) are held fixed at their \texttt{astromodels} defaults, so
the free parameters are the indices, breaks and normalization. The blackbody is
$N(E)\propto E^2/[\exp(E/\kT)-1]$. For all peaked models \Ep\ is the true
$\nu F_\nu$ peak. The $+$BB models and the 2SBPL are the competing descriptions of
curvature beyond a single break --- a thermal bump versus a second synchrotron break.

\subsection{Model selection} \label{sec:select}
We rank models by the Akaike and Bayesian information criteria
\citep[AIC, BIC;][]{Akaike1974,Schwarz1978} using the standard $-2\ln L$, and by
nested likelihood-ratio tests \citep[LRT;][]{Wilks1938} for Band$+$BB/Band,
CPL$+$BB/CPL, and 2SBPL/SBPL; a blackbody or second break is significant at
$\mathrm{LRT}\ge9.2$ (99\%, 2 d.o.f.). A physical-validity gate prevents a fit
from \emph{winning} if any key parameter is railed against a bound or, for the
2SBPL, if the breaks are mis-ordered ($x_b\ge x_p$). A bin counts as
``two-break'' or ``thermal-proxy'' only if the corresponding curved model is
physically valid.

\subsection{Robustness of the blackbody fit} \label{sec:multistart}
The blackbody has a railing local minimum ($\kT\!\to\!1$~keV, normalization
$\to0$) that a poor seed can fall into; since per-bin fits are seeded from the
time-integrated fit, a railed seed can propagate to every bin. We therefore
multi-start each $+$BB model from a grid of temperature seeds and retain the
lowest $-2\ln L$ (e.g.\ for GRB~110721A this reduces railed bins from 12/13 to
0/13). The full-sample fits used for the sample-wide BB statistics predate this
fix and therefore \emph{under-detect} the blackbody; we treat the resulting
thermal fractions as conservative (\S\ref{sec:curv}). The 2SBPL/break results do
not involve the blackbody and are unaffected.

\subsection{Correlation fitting} \label{sec:dagostini}
For every correlation we report the Spearman rank coefficient $\rho$ and its
chance probability $p$, and we fit a power law in log--log space with the
\citet{DAgostini2005} errors-in-both-variables likelihood, which simultaneously
infers the slope $m$, normalization $K$ at a fixed pivot, and the
\emph{intrinsic scatter} \ssc\ (dex), marginalized by MCMC. We report \ssc\
rather than a reduced $\chi^2$, since it directly quantifies the dispersion not
accounted for by measurement error.

% ============================================================================
\section{Results} \label{sec:results}
% ============================================================================

\subsection{Low-energy index and the model mix} \label{sec:coverage}
All 106 bursts were fit (1063 bins $\times$ 6 models). The low-energy photon index
peaks at a median Band $\alpha\simeq-0.75$ (Figure~\ref{fig:lowE}), between the
slow-cooled ($-2/3$) and fast-cooled ($-3/2$) synchrotron limits and close to the
$-0.81\pm0.1$ found by \citet{Burgess2014b}; the population sits near the
slow-cooling line, with relatively few bins as soft as $-3/2$. Among single-break
models the lowest-AIC winner is split between CPL, SBPL and Band, reflecting the
modest per-bin signal-to-noise of single-pulse bins.

\begin{figure}[t]
\centering
\includegraphics[width=\columnwidth]{figures/fig_lowE_index.png}
\caption{Distribution of the low-energy photon index (Band $\alpha$, CPL index,
SBPL $\alpha$) across the 1063 bins. Dashed/dotted lines mark the slow- ($-2/3$)
and fast-cooled ($-3/2$) synchrotron limits; the population peaks near $-0.75$.}
\label{fig:lowE}
\end{figure}

\subsection{Curvature: thermal proxy versus a second break} \label{sec:curv}
Restricting to bins requiring curvature beyond a single break ($\daic>6$ relative
to the best single-break model, with the curved model physically valid), we
classify each as thermal-proxy (BB$+$SBPL), two-break (2SBPL) or degenerate
(Table~\ref{tab:curv}; Figure~\ref{fig:curv}). At $\daic>6$, of 82 curved bins
$78\%$ are thermal-or-degenerate and $22\%$ genuinely prefer the explicit second
break; at $\daic\ge10$ the split is $82\%$/$18\%$. Moreover $94\%$ of all 1063
bins cannot separate their best two models at $\daic\ge10$ --- single-pulse GBM
bins rarely have the S/N to decide among six models. The thermal flux fraction in
BB-significant bins is modest, $F_{\rm BB}/F_{\rm tot}\simeq0.11$
(16--84\%: 0.05--0.24), below the $0.27$--$0.39$ of \citet{Burgess2014a};
because the empirical continuum can absorb thermal flux and the present fits
under-detect the BB (\S\ref{sec:multistart}), these thermal fractions are
conservative. The brightest burst, GRB~130427A, provides the strongest genuine
two-break detections.

\begin{figure}[t]
\centering
\includegraphics[width=\columnwidth]{figures/fig_curvature_split.png}
\caption{Of bins requiring curvature beyond a single break, the split between a
thermal proxy (BB$+$SBPL), a genuine second break (2SBPL), and degenerate cases.}
\label{fig:curv}
\end{figure}

\subsection{A worked example: GRB~130427A} \label{sec:130427}
Figure~\ref{fig:evol} shows the time-resolved $\nu F_\nu$ evolution of GRB~130427A
(88 bins), the brightest and best-resolved burst in the sample, color-coded by
time. The spectrum hardens to the pulse peak and softens through the decay
(hard-to-soft evolution), and this burst alone yields the strongest 2SBPL
detections (peak-bin $\daic\sim590$ over the best single-break model) and a clean
$\Ep$--$\kT$ track (\S\ref{sec:burgess}).

\begin{figure}[t]
\centering
\includegraphics[width=\columnwidth]{figures/fig_nufnu_evolution_130427A.png}
\caption{Time-resolved $\nu F_\nu$ evolution of GRB~130427A (best-fit model per
bin, peak-normalized, colored by time). Hard-to-soft spectral evolution is clear.}
\label{fig:evol}
\end{figure}

\subsection{The Burgess $\Ep$--$\kT$ correlation in a large sample}
\label{sec:burgess}
We test the \citet{Burgess2014a} thermal--non-thermal correlation per burst, using
\Ep\ and \kT\ from $+$BB fits in BB-significant bins. Of 24 bursts with $\ge3$ such
bins, only \textbf{12 have $\ge5$} (Table~\ref{tab:perburst}); the slope is robust
only for the highest-S/N bursts, led by GRB~130427A ($N=36$, $\Ep\propto\kT^{0.66}$,
$\rho=0.78$, consistent with the published $\alpha=1.02$). Mapping the slope to jet
composition in the \citeauthor{Burgess2014a} scheme ($\alpha\!\lesssim\!1.5$
baryon-dominated, $\gtrsim\!1.5$ magnetically-dominated), the well-constrained
bursts are predominantly baryonic. The fainter bursts yield too few
BB-significant bins for a reliable per-burst slope (Figure~\ref{fig:burgess});
this is a consequence of \emph{empirical} modeling --- the flexible Band/CPL
continuum absorbs thermal curvature, so an empirical $+$BB detects a blackbody in
fewer bins than the physical synchrotron$+$BB model of \citeauthor{Burgess2014a}
(quantified by the low $F_{\rm BB}/F_{\rm tot}$, \S\ref{sec:curv}). The
correlation is thus present where the S/N permits its detection, but the empirical
approach recovers it less completely than a physical model.

\begin{figure}[t]
\centering
\includegraphics[width=\columnwidth]{figures/fig_burgess_reblock.png}
\caption{Per-burst $\Ep$--$\kT$ for the six \citet{Burgess2014a} bursts (black:
BB-significant bins; grey: production bins). Only GRB~130427A (bottom-right) has
enough significant bins for a reliable slope.}
\label{fig:burgess}
\end{figure}

\subsection{A correlation between the two spectral breaks} \label{sec:nucnum}
We now turn to a property not previously quantified: the relation between the two
2SBPL break energies. Across all bins with a valid, significant 2SBPL
($\mathrm{LRT_{2SBPL/SBPL}}\ge9.2$, breaks non-railed and ordered; $N=40$ from 22
bursts), the injection and cooling breaks are correlated,
\begin{equation}
\num \;=\; K\left(\frac{\nuc}{100~{\rm keV}}\right)^{m}, \quad
m=0.73\pm0.10,\ \ssc=0.26~{\rm dex},
\end{equation}
with Spearman $\rho=0.74$, $p=6\times10^{-8}$ (Table~\ref{tab:corr};
Figure~\ref{fig:corr}A). Because the validity gate enforces $x_b<x_p$, we verified
by permutation that this is not a truncation artifact: the observed $\rho=0.74$
exceeds the 95th percentile of the truncated null ($\rho_{95}=0.54$). The thermal
proxy reproduces the same trend: the analogous pair correlates as
$\Ep\propto\kT^{0.60\pm0.08}$ ($\rho=0.54$, $\ssc=0.50$~dex, $N=198$;
Figure~\ref{fig:corr}B), so the result is robust to the choice of curvature model.
To our knowledge this is the first quantified correlation between the two
synchrotron break energies: \citet{Mei2024} correlate each break with luminosity
separately and use $\num/\nuc$ only as a cooling-regime classifier, while
\citet{Ravasio2019} display the two breaks together (their Fig.~4) but report no
fitted relation.

\begin{figure*}[t]
\centering
\includegraphics[width=\textwidth]{figures/fig_break_correlations.png}
\caption{Spectral-break/temperature correlations (full sample, observer frame).
(A) The two 2SBPL breaks $\num$--$\nuc$. (B) The BB$+$SBPL proxy $\Ep$--$\kT$.
(C) $\nuc$ (Band \Ep) and (D) \kT\ versus observer-frame energy flux; (E) the
$\num/\nuc$ cooling-regime distribution. Panels C--D are annotated with the
induction-controlled (photon-flux) Spearman coefficient (\S\ref{sec:lum}).}
\label{fig:corr}
\end{figure*}

\subsection{Break/temperature energy versus brightness} \label{sec:lum}
Motivated by \citet{Mei2024} --- who find a rest-frame $\nuc$--\Liso\ relation but
no $\Ep$--\Liso\ relation --- we test the observer-frame analog. Our sample lacks
redshifts, so we use the per-bin observer-frame energy flux (10~keV--40~MeV) as a
luminosity proxy; this is \emph{not} rest-frame \Liso\ and we interpret it
cautiously. Energy flux is inflated by spectral hardness, which can manufacture a
peak-energy--flux correlation; we control for this by repeating each test against
the \emph{photon} flux (no energy weighting). After this control
(Table~\ref{tab:corr}; Figure~\ref{fig:corr}C--D): the $\Ep$--flux (Yonetoku-like)
correlation \emph{collapses} from $\rho=0.38$ to $\rho=0.01$ --- it is an
energy-weighting artifact, consistent with the absence of an $\Ep$--\Liso\
relation reported by \citeauthor{Mei2024}; the $\nuc$--flux correlation is weak
($\rho=0.51\!\to\!0.19$); and the blackbody-temperature--brightness correlation
\emph{survives}, $\rho=0.60\!\to\!0.35$ ($p=5\times10^{-7}$), indicating a genuine
$\kT$--brightness link.

% ============================================================================
\section{Discussion} \label{sec:discussion}
% ============================================================================
\subsection{Empirical versus physical modeling}
Our analysis is deliberately empirical: Band/CPL/SBPL/2SBPL/$+$BB are
model-agnostic and applied uniformly, which is ideal for population statistics and
for the robust \emph{detection} of correlations across a large $N$. The cost is
that a flexible empirical continuum can absorb a real sub-dominant component ---
quantified here by the low $F_{\rm BB}/F_{\rm tot}\simeq0.1$ and the recovery of
the $\Ep$--$\kT$ correlation only at high S/N. This inverts the
\citet{Burgess2014b} argument: where they show a physical synchrotron model
separates components more cleanly than empirical fitting, we show that a large
empirical survey is nonetheless sufficient to \emph{detect} the key correlations,
while confirming that the per-burst thermal--non-thermal relation at their fidelity
requires the physical model. We do not fit a physical synchrotron model; every
physical mapping below is an interpretation of the empirical observables, not a
direct inference.

\subsection{Physical interpretation}
The low-energy index peaking at $\alpha\simeq-0.75$, between $-2/3$ and $-3/2$ and
near the slow-cooling line, is consistent with marginally-fast/slow-cooled
synchrotron with ongoing acceleration, as \citet{Burgess2014b} concluded. For the
break correlation, standard synchrotron scalings give $\num\propto B\gamma_m^2$ and
$\nuc\propto B^{-3}t^{-2}$ \citep{Sari1998}; the observed sub-linear
$\num\propto\nuc^{0.73}$ then implies that the injection break rises more slowly,
relative to the cooling break, than a fixed magnetic-field/geometry would predict
--- i.e.\ the spectral separation $\num/\nuc$ is not constant but evolves, a
constraint on the joint evolution of $B$, $\gamma_m$ and $\gamma_c$. For the
thermal picture, mapping $\Ep\propto\kT^{\alpha}$ to the outflow via
$\Gamma(r)\propto r^{\mu}$ \citep{Burgess2014a}, $\alpha\!\sim\!1$ ($\mu\!\sim\!1$)
indicates a baryon-dominated (kinetic) jet and $\alpha\!\sim\!2$ a
magnetically-dominated jet; our well-constrained bursts (Table~\ref{tab:perburst})
are predominantly baryonic, although the per-burst slopes are uncertain.

\subsection{Comparison with \citet{Mei2024}}
Our induction-controlled flux analysis independently reproduces the central
\citeauthor{Mei2024} result: the apparent $\Ep$--luminosity (Yonetoku) correlation
is not fundamental (it vanishes against photon flux), while a genuine
$\kT$--brightness link points to a thermal--energetics relation. A definitive
comparison to the rest-frame $\nuc$--\Liso\ relation requires redshifts, which
our sample lacks.

\subsection{Limitations and next steps}
The sample-wide BB statistics use pre-multi-start fits and under-detect the
blackbody; a full re-run with the multi-start engine (\S\ref{sec:multistart}) will
firm up the thermal fractions and per-burst slopes (the curvature split and
$\kT$ correlations may strengthen). A redshift cross-match of the bright subset
would enable the rest-frame $\nuc$--\Liso\ test. Finally, the time-resolved bins
invite a hardness-intensity / hardness-fluence and a $\kT(t)$, $R(t)$ evolution
analysis \citep{Burgess2014b}, which we defer to future work.

% ============================================================================
\section{Conclusions} \label{sec:conclusions}
% ============================================================================
\begin{enumerate}
\item We performed time-resolved \emph{empirical} spectroscopy of \textbf{106
single-pulse bright Fermi/GBM bursts} (1063 Bayesian-block bins; six models each)
--- among the largest uniform single-pulse spectral surveys to date.
\item The low-energy index peaks at $\alpha\simeq-0.75$, near the slow-cooled
synchrotron limit and close to \citeauthor{Burgess2014b}'s $-0.81$.
\item Of bins requiring curvature beyond a single break, $78\%$ ($\daic>6$;
$82\%$ at $\daic\ge10$) are equally or better described by a thermal proxy than by
an explicit second break; $94\%$ of all bins are statistically inconclusive among
the six models; $F_{\rm BB}/F_{\rm tot}\simeq0.1$.
\item The \citet{Burgess2014a} $\Ep$--$\kT$ correlation is recovered for the
highest-S/N bursts (GRB~130427A, $\rho=0.78$, baryonic) but not robustly in
fainter bursts with empirical models --- a model-sensitivity, not a sample, limit.
\item We report a \textbf{new correlation between the two 2SBPL break energies},
$\num\propto\nuc^{0.73\pm0.10}$ ($\rho=0.74$, $\ssc=0.26$~dex), mirrored by the
thermal proxy, and not previously quantified.
\item In an observer-frame proxy we find no genuine $\Ep$--brightness correlation
(consistent with \citealt{Mei2024}) but a real $\kT$--brightness correlation.
\end{enumerate}

% ============================================================================
% TABLES
% ============================================================================
\begin{deluxetable}{lccccc}
\tablecaption{Sample (representative rows; full table machine-readable).\label{tab:sample}}
\tablehead{\colhead{GRB} & \colhead{$T_{90}$} & \colhead{Fluence} & \colhead{NaI/BGO} & \colhead{LLE} & \colhead{$N_{\rm bins}$}\\
\colhead{} & \colhead{(s)} & \colhead{(erg\,cm$^{-2}$)} & \colhead{} & \colhead{} & \colhead{}}
\startdata
130427A & 13.6 & $2.5\times10^{-3}$ & 3/1 & Y & 88 \\
110721A & 21.8 & $3.5\times10^{-5}$ & 4/1 & Y & 13 \\
150902733 & 13.6 & $1.6\times10^{-4}$ & 3/1 & Y & 16 \\
\nodata & \nodata & \nodata & \nodata & \nodata & \nodata \\
\enddata
\tablecomments{106 bursts total (\citealt{Busby2024} single-pulse selection,
score $\gtrsim0.9983$). The full machine-readable table lists all 106 bursts.}
\end{deluxetable}

\begin{deluxetable*}{lcccccc}
\tablecaption{Correlation results (full sample, observer frame). Slopes from the
\citet{DAgostini2005} intrinsic-scatter fit.\label{tab:corr}}
\tablehead{\colhead{Relation ($y$ vs $x$)} & \colhead{$N$} & \colhead{$\rho$} & \colhead{$p$} & \colhead{slope $m$} & \colhead{\ssc\ (dex)} & \colhead{note}}
\startdata
$\num$ vs $\nuc$ (2SBPL)        & 40  & 0.74 & $6\times10^{-8}$  & $0.73\pm0.10$ & 0.26 & null $\rho_{95}=0.54$\\
$\Ep$ vs $\kT$ (BB proxy)       & 198 & 0.54 & $1\times10^{-16}$ & $0.60\pm0.08$ & 0.50 & \\
$\nuc$ vs energy flux           & 933 & 0.51 & $1\times10^{-63}$ & $0.88\pm0.05$ & \nodata & artifact-prone\\
$\nuc$ vs photon flux           & 933 & 0.19 & $4\times10^{-9}$  & \nodata       & \nodata & induction control\\
$\Ep$ vs energy flux            & 198 & 0.38 & $2\times10^{-8}$  & \nodata       & \nodata & artifact\\
$\Ep$ vs photon flux            & 198 & 0.01 & 0.9               & \nodata       & \nodata & $\Rightarrow$ no real corr.\\
$\kT$ vs energy flux            & 198 & 0.60 & $8\times10^{-21}$ & \nodata       & \nodata & \\
$\kT$ vs photon flux            & 198 & 0.35 & $5\times10^{-7}$  & \nodata       & \nodata & real residual\\
\enddata
\tablecomments{Energy flux (10\,keV--40\,MeV, observer frame) is a luminosity
proxy only; no redshifts. The energy-/photon-flux pairs isolate the
energy-weighting artifact.}
\end{deluxetable*}

\begin{deluxetable}{lccc}
\tablecaption{Per-burst $\Ep$--$\kT$ slopes ($N\ge5$ BB-significant bins).\label{tab:perburst}}
\tablehead{\colhead{GRB} & \colhead{$N$} & \colhead{$\alpha$ ($\Ep\propto\kT^\alpha$)} & \colhead{jet}}
\startdata
130427A & 36 & 0.66 & baryonic \\
230812790 & 14 & 1.27 & baryonic \\
160910722 & 8 & 0.75 & baryonic \\
100707032 & 7 & 1.51 & magnetic \\
170921168 & 7 & 1.10 & baryonic \\
210714331 & 7 & 0.60 & baryonic \\
191125206 & 6 & 0.77 & baryonic \\
221209243 & 6 & 1.67 & magnetic \\
231020790 & 6 & 0.92 & baryonic \\
250313607 & 5 & 0.53 & baryonic \\
250407659 & 5 & 1.41 & baryonic \\
180723757 & 5 & $-1.32$ & \nodata \\
\enddata
\tablecomments{Jet type from $\alpha\!<\!1.5$ (baryonic) / $\ge1.5$ (magnetic),
\citep{Burgess2014a}. Slopes are uncertain at low $N$; 180723757 is
scatter-dominated.}
\end{deluxetable}

\begin{deluxetable}{lcccc}
\tablecaption{Curvature classification (physically-valid curved models).\label{tab:curv}}
\tablehead{\colhead{Threshold} & \colhead{curved} & \colhead{thermal-proxy} & \colhead{degenerate} & \colhead{two-break}}
\startdata
$\daic>6$    & 82 & 55 & 9 & 18 \\
$\daic\ge10$ & 51 & 33 & 9 & 9 \\
\enddata
\tablecomments{$78\%$ ($\daic>6$) / $82\%$ ($\daic\ge10$) thermal-or-degenerate.
$94\%$ of all 1063 bins are inconclusive (best two models within $\daic<10$).
Thermal counts are conservative (BB under-detected; \S\ref{sec:multistart}).}
\end{deluxetable}

\begin{acknowledgments}
% %TODO grants, facility/data acknowledgments.
This work made use of \textsf{3ML} \citep{Vianello2015} and the Fermi/GBM and LAT
public data. Bayesian blocks follow \citet{Scargle2013}; correlation fits use the
\citet{DAgostini2005} method.
\end{acknowledgments}

\facilities{Fermi(GBM), Fermi(LAT)}
\software{3ML \citep{Vianello2015}, Astropy, Matplotlib, NumPy, SciPy, emcee}

% ============================================================================
% Author-year natbib: EVERY \bibitem optional arg MUST contain (YYYY).
% %chk cross-check all volume/page numbers against ADS before submission.
% ============================================================================
\begin{thebibliography}{}
\bibitem[Akaike(1974)]{Akaike1974} Akaike, H. 1974, IEEE Trans. Autom. Control, 19, 716
\bibitem[Atwood et al.(2009)]{Atwood2009} Atwood, W.~B., Abdo, A.~A., Ackermann, M., et al. 2009, \apj, 697, 1071
\bibitem[Band et al.(1993)]{Band1993} Band, D., Matteson, J., Ford, L., et al. 1993, \apj, 413, 281
\bibitem[Burgess et al.(2014a)]{Burgess2014a} Burgess, J.~M., Preece, R.~D., Connaughton, V., et al. 2014a, \apjl, 784, L43
\bibitem[Burgess et al.(2014b)]{Burgess2014b} Burgess, J.~M., Preece, R.~D., Connaughton, V., et al. 2014b, \apj, 784, 17
\bibitem[Busby \& Lazzati(2024)]{Busby2024} Busby, K., \& Lazzati, D. 2024, \apj, 972, 83
\bibitem[D'Agostini(2005)]{DAgostini2005} D'Agostini, G. 2005, arXiv:physics/0511182
\bibitem[Goldstein et al.(2012)]{Goldstein2012} Goldstein, A., Burgess, J.~M., Preece, R.~D., et al. 2012, \apjs, 199, 19
\bibitem[Gruber et al.(2014)]{Gruber2014} Gruber, D., Goldstein, A., Weller von Ahlefeld, V., et al. 2014, \apjs, 211, 12
\bibitem[Guiriec et al.(2011)]{Guiriec2011} Guiriec, S., Connaughton, V., Briggs, M.~S., et al. 2011, \apjl, 727, L33
\bibitem[Kaneko et al.(2006)]{Kaneko2006} Kaneko, Y., Preece, R.~D., Briggs, M.~S., et al. 2006, \apjs, 166, 298
\bibitem[Kumar \& Zhang(2015)]{KumarZhang2015} Kumar, P., \& Zhang, B. 2015, \physrep, 561, 1
\bibitem[Meegan et al.(2009)]{Meegan2009} Meegan, C., Lichti, G., Bhat, P.~N., et al. 2009, \apj, 702, 791
\bibitem[Mei et al.(2025)]{Mei2024} Mei, A., Oganesyan, G., \& Macera, S. 2025, \aap, 693, A156
\bibitem[M\'esz\'aros \& Rees(2000)]{MeszarosRees2000} M\'esz\'aros, P., \& Rees, M.~J. 2000, \apj, 530, 292
\bibitem[Oganesyan et al.(2017)]{Oganesyan2017} Oganesyan, G., Nava, L., Ghirlanda, G., \& Celotti, A. 2017, \apj, 846, 137
\bibitem[Oganesyan et al.(2018)]{Oganesyan2018} Oganesyan, G., Nava, L., Ghirlanda, G., \& Celotti, A. 2018, \aap, 616, A138
\bibitem[Oganesyan et al.(2019)]{Oganesyan2019} Oganesyan, G., Nava, L., Ghirlanda, G., Melandri, A., \& Celotti, A. 2019, \aap, 628, A59
\bibitem[Pe'er et al.(2007)]{Peer2007} Pe'er, A., Ryde, F., Wijers, R.~A.~M.~J., M\'esz\'aros, P., \& Rees, M.~J. 2007, \apjl, 664, L1
\bibitem[Preece et al.(2014)]{Preece2014} Preece, R., Burgess, J.~M., von Kienlin, A., et al. 2014, Science, 343, 51
\bibitem[Ravasio et al.(2018)]{Ravasio2018} Ravasio, M.~E., Oganesyan, G., Ghirlanda, G., et al. 2018, \aap, 613, A16
\bibitem[Ravasio et al.(2019)]{Ravasio2019} Ravasio, M.~E., Ghirlanda, G., Nava, L., \& Ghisellini, G. 2019, \aap, 625, A60
\bibitem[Ryde(2004)]{Ryde2004} Ryde, F. 2004, \apj, 614, 827
\bibitem[Ryde \& Pe'er(2009)]{RydePeer2009} Ryde, F., \& Pe'er, A. 2009, \apj, 702, 1211
\bibitem[Sari et al.(1998)]{Sari1998} Sari, R., Piran, T., \& Narayan, R. 1998, \apjl, 497, L17
\bibitem[Scargle et al.(2013)]{Scargle2013} Scargle, J.~D., Norris, J.~P., Jackson, B., \& Chiang, J. 2013, \apj, 764, 167
\bibitem[Schwarz(1978)]{Schwarz1978} Schwarz, G. 1978, Ann. Stat., 6, 461
\bibitem[Vianello et al.(2015)]{Vianello2015} Vianello, G., Lauer, R.~J., Younk, P., et al. 2015, arXiv:1507.08343
\bibitem[Wilks(1938)]{Wilks1938} Wilks, S.~S. 1938, Ann. Math. Stat., 9, 60
\bibitem[Yonetoku et al.(2004)]{Yonetoku2004} Yonetoku, D., Murakami, T., Nakamura, T., et al. 2004, \apj, 609, 935
\bibitem[Yu et al.(2016)]{Yu2016} Yu, H.-F., Preece, R.~D., Greiner, J., et al. 2016, \aap, 588, A135
\end{thebibliography}

\end{document}
